%% LyX 2.4.4 created this file.  For more info, see https://www.lyx.org/.
%% Do not edit unless you really know what you are doing.
\documentclass[english]{article}
\usepackage[T1]{fontenc}
\usepackage[latin9]{inputenc}
\usepackage{enumitem}
\usepackage{amsmath}
\usepackage{amsthm}
\usepackage{cancel}
\usepackage{geometry}
\geometry{verbose,tmargin=1in,bmargin=1in,lmargin=1in,rmargin=1in}

\makeatletter
%%%%%%%%%%%%%%%%%%%%%%%%%%%%%% Textclass specific LaTeX commands.
\numberwithin{equation}{section}
\numberwithin{figure}{section}
\newlength{\lyxlabelwidth}      % auxiliary length 

%%%%%%%%%%%%%%%%%%%%%%%%%%%%%% User specified LaTeX commands.
\usepackage{fancyhdr}
\usepackage{lastpage}
\pagestyle{fancy}
\usepackage{enumitem}
\usepackage{ifthen}
%sets math fonts to be more readable
\everymath=\expandafter{\the\everymath\displaystyle}
%\DeclareMathSizes{10}{10}{10}{10}
\usepackage{multicol}

\usepackage{tikz}
\usetikzlibrary{plotmarks}
\usepackage{pgfplots}
\pgfplotsset{compat=newest}

\usepackage{calc}
\usepackage[nomessages]{fp}

\usepackage{advdate}    % Advancing/saving dates
\usepackage{datetime}   % Dates formatting
\usepackage{datenumber} % Counters for dates
% Added by lyx2lyx
\setlength{\parskip}{\smallskipamount}
\setlength{\parindent}{0pt}

\makeatother

\usepackage{babel}
\begin{document}
\renewcommand{\author}{Dr.\,James Ashe}

\newcommand{\classNum}{132}

\newcommand{\classSec}{C \& K}

\newcommand{\nameType}{Name:}

\newcommand{\examType}{Worksheet 1.1 \& 1.2}

\renewcommand{\date}{\formatdate{14}{1}{2014}}

%%First page's header%%%%%%%%%%%%%%%%%%%%%%
\fancypagestyle{firststyle} %%header settings for first pagr
{
	\fancyhead[L]{MTH \classNum{}(\classSec{}) \author{}\\
	\textbf{\examType{}} \date{}}
	\fancyhead[C]{\textbf{\nameType}}
	\setlength{\headsep}{14pt}
}
%%%%%%%%%%%%%%%%%%%%%%%%%%%%%%%%%%%%%%%%%%

%%Next page's header%%%%%%%%%%%%%%%%%%%%%%%
\setlist{nolistsep}
\lhead{MTH \classNum{}(\classSec{})} 
\chead{\textbf{\nameType}} 
\cfoot{\thepage{}~of~\pageref{LastPage}} 
\setlength{\headsep}{5pt} 
%%%%%%%%%%%%%%%%%%%%%%%%%%%%%%%%%%%%%%%%%%%

%%Various settings; see the preamble for other settings%%%%%
%%\setenumerate[0]{leftmargin=12pt,itemindent=0pt,labelwidth=0pt}
\renewcommand{\labelenumi}{\textbf{\arabic{enumi}.}}
\renewcommand{\labelenumii}{\textbf{\alph{enumii}.}}
\thispagestyle{firststyle}
%%%%%%%%%%%%%%%%%%%%%%%%%%%%%%%%%%%%%%%%%%
\newcommand{\xmax}{10}
\newcommand{\xmin}{-10}
\newcommand{\xstep}{0} %must be xmin+n
\newcommand{\ymax}{10}
\newcommand{\ymin}{-10}
\newcommand{\ystep}{0} %must be ymin+n

\newcommand{\xspace}{\vspace{1.2cm}}

\textbf{Use the given conditions to write an equation for the line
in slope-intercept form, if possible.}
\begin{enumerate}[resume]
\item Slope$=-3$, passing through $(2,6)$.\vfill
\item The line is parallel to $y=-\frac{1}{2}x+5$ and passes through $(2,0)$.\vfill
\item Through $(5,0)$ and $(-6,2)$.\vfill
\item Through $(1,2)$ and $(3,2)$.\vfill
\item Through $(3,1)$ and $(3,2)$.\vfill\newpage
\end{enumerate}
\textbf{Solve the problem.}
\begin{enumerate}[resume]
\item \label{enu:After-5-years}After 5 years as an actuary, Sue's salary
was \$130,000. After 12 years her salary had risen to \$200,000. Assuming
that the change in her salary over time is linear, find an equation
expressing her salary as $y$ and $x$ as the number of years on the
job.\vfill
\item Using the equation you found in \ref{enu:After-5-years}, how many
years does she need to work as an actuary to have a salary of \$250,000?\vfill
\item A shoe company will make a new type of shoe which costs \$21 to make
per shoe. The fixed cost of producing this shoe (utilities, property
taxes, leases, etc.) is \$25,000. The company plans to sell the shoe
for \$100 a pair. 

\begin{enumerate}[resume]
\item Find the cost function in dollars, $C(x)$, for $x$ pairs of shoes.
\xspace
\item Find the revenue function in dollars, $R(x)$, for $x$ pairs of shoes.\xspace
\item Find the profit function in dollars, $P(x)$, for $x$ pairs of shoes.\xspace\xspace
\item How many pair of shoes must they sell to break even?\vfill
\end{enumerate}
\end{enumerate}

\end{document}
